\chapter{Solutions}

\section{Exercices essentiels}\label{I-exercices-essentiels}

\runinhead{Propriétés
élémentaires}\phantomsection\addcontentsline{toc}{subsubsection}{Propriétés
élémentaires}\label{I-answer-propelem}

\begin{enumerate}
\def\labelenumi{\arabic{enumi}.}
\item
  \(\mathbb{P}(A) \in [0,1]\) est donné par la définition. Par ailleurs,
  on a \(\Omega = A \cup A^c\) et \(A\cap A^c = \varnothing\). En
  considérant la suite \(A_1 = A\), \(A_2 = A^c\) et
  \(A_n = \varnothing\) pour \(n > 2\), le point 2. de la définition
  nous donne
  \(1 = \mathbb{P}(\Omega) = \mathbb{P}(A) + \mathbb{P}(A^c)\).
\item
  Puisque \(A \subset B\), on a l'union disjointe
  \(B = (B\cap A) \cup (B\cap A^c) = A \cup (B\cap A^c)\). Comme
  \(\mathbb{P}(B\cap A^c) \geq 0\) par définition, on en déduit le
  résultat.
\item
  On a \((A\cup B) = B \cup (A\cap B^c)\), avec
  \(B \cap (A\cap B^c) =\varnothing\) ainsi que
  \(A = (A\cap B^c) \cup (A\cap B)\), où
  \((A\cap B^c) \cap (A\cap B) = \varnothing\). On en déduit
  \(\mathbb{P}(A\cup B) = \mathbb{P}(B) + \mathbb{P}(A\cap B^c)\) et
  \(\mathbb{P}(A) = \mathbb{P}(A\cap B^c) + \mathbb{P}(A\cap B)\), d'où
  le résultat.
\item
  Par récurrence. On a déjà vu que
  \(\mathbb{P}(A\cup B) \leq \mathbb{P}(A) + \mathbb{P}(B)\). Notant
  \(B = \bigcup_{i=1}^n A_i\), alors
  \(\mathbb{P}(\bigcup_{i=1}^{n+1} A_i) \leq \mathbb{P}(B) + \mathbb{P}(A_{n+1})\),
  et l'hypothèse de récurrence nous indique
  \(\mathbb{P}\left(\bigcup_{i=1}^n A_i\right) \leq \sum_{i=1}^n \mathbb{P}(A_i)\),
  d'où le résultat.
\item
  Par récurrence. On a déjà vu que
  \(\mathbb{P}(A \cup B ) = \mathbb{P}(A) + \mathbb{P}(B) - \mathbb{P}(A \cap B)\).
  Notant \(B = \bigcup_{i=1}^n A_i\), alors \begin{align*}
      \mathbb{P}(\bigcup_{i=1}^{n+1} A_i) =& \mathbb{P}(B \cup A_{n+1}) \\
                                  =& \mathbb{P}(B) + \mathbb{P}(A_{n+1}) - \mathbb{P}(B \cap A_{n+1}) \\
                                  =& \sum_{i=1}^n \mathbb{P}(A_i) - \sum_{1 \leq i < j \leq n} \mathbb{P}(A_i \cap A_j) + \ldots + (-1)^n \mathbb{P}\left(\bigcap_{i=1}^n A_i\right) \\
                                   &+ \mathbb{P}(A_{n+1})\\
                                   &-\sum_{i=1}^n \mathbb{P}(A_i\cap A_{n+1}) + \sum_{1 \leq i < j \leq n} \mathbb{P}(A_i \cap A_j) + \ldots + (-1)^{n+1} \mathbb{P}\left(\bigcap_{i=1}^n A_i\cap A_{n+1}\right)\\
                                   =& \sum_{i=1}^{n+1} \mathbb{P}(A_i) - \sum_{1 \leq i < j \leq n +1 } \mathbb{P}(A_i \cap A_j) + \ldots + (-1)^{n+1} \mathbb{P}\left(\bigcap_{i=1}^{n+1} A_i\right).
  \end{align*}
\end{enumerate}

\runinhead{Continuité
monotone}\phantomsection\addcontentsline{toc}{subsubsection}{Continuité
monotone}\label{I-answer-contmon}
On définit une suite \((B_n)_{n\in \mathbb{N}^\ast}\) telle que
\(B_1 = A_1\) et \(B_n = A_n \setminus A_{n-1}\), pour \(n\geq 2\). Les
\(B_n\) ainsi définis sont deux à deux disjoints et
\(\cup B_n = \cup A_n\), donc~:
\[ \mathbb{P}\left(\bigcup_{n\in\mathbb{N}^\ast} A_n\right) = \sum_n \mathbb{P}(B_n) = \lim_{n \to \infty} \sum_{p=1}^n \mathbb{P}(B_p) = \lim_{n \to \infty} \mathbb{P}(A_n).\]

\runinhead{Une définition alternative de la
probabilité}\phantomsection\addcontentsline{toc}{subsubsection}{Une
définition alternative de la probabilité}\label{I-answer-altdef}
Soit \(A_n\) une suite d'éléments de \(\mathcal{A}\) deux-à-deux
disjoints. On définit \(B_n = \cup_{p \leq n} A_n\) et
\(B = \cup_n A_n\). Comme \(\mathbb{P}\) est additive, on a
\(\mathbb{P}(B_n) = \sum_{p \leq n} \mathbb{P}(A_n)\) qui croît vers
\(\sum_n \mathbb{P}(A_n)\) et aussi vers \(\mathbb{P}(B)\) d'après le
point 3.

En considérant les résultats de ces deux exercices, on obtient une
définition alternative de la probabilité en substituant la
\protect\hyperref[I-contmon]{continuité
monotone (p.~\pageref*{I-contmon})} et l'additivité à la propriété de
\protect\hyperref[I-defproba]{\(\sigma\)-additivité (p.~\pageref*{I-defproba})}.

\runinhead{Inégalité}\phantomsection\addcontentsline{toc}{subsubsection}{Inégalité}\label{I-answer-boolinf}
On a déjà vu le cas fini avec \protect\hyperref[I-propelem]{l'inégalité
de Boole (p.~\pageref*{I-propelem})}. On pose
\(B_n = \cup_{i=1}^n A_i\), qui croît vers l'ensemble
\(C = \cup_n A_n\). D'après \protect\hyperref[I-propelem]{l'inégalité de
Boole (p.~\pageref*{I-propelem})}, on a
\[\mathbb{P}(B_n) \leq \sum_{i=1}^n \mathbb{P}(A_i)\] Mais
\(\mathbb{P}(B_n) \to_{n \to \infty} \mathbb{P}(C)\) d'après le
\protect\hyperref[I-contmon]{théorème de continuité
monotone (p.~\pageref*{I-contmon})}, tandis que
\(\sum_{i=1}^n \mathbb{P}(A_i) \to_{n \to \infty} \sum_{n\in\mathbb{N}^\ast} \mathbb{P}(A_n)\).
En passant à la limite, on obtient donc le résultat.

\runinhead{Test}\phantomsection\addcontentsline{toc}{subsubsection}{Test}\label{I-answer-test}
On considère les événements \(A\) ``l'individu est séropositif'', et
\(B\) ``le test de détection donne un résultat positif''. Les données
ci-dessus nous indiquent \(P(A) = 10^{-4}\) d'où
\(\mathbb{P}(A^c) = 0,9999\) , \(\mathbb{P}(B|A) = 0,99\) et
\(\mathbb{P}(B|A^c) = 0,001\). Nous trouvons alors \begin{align*}
\mathbb{P}(A|B) &= \frac{\mathbb{P}(A \cap B)}{\mathbb{P}(B)}\\
        &= \frac{\mathbb{P}(B |A ) \mathbb{P}(A)}{\mathbb{P}(B |A)\mathbb{P}(A) + \mathbb{P}(B|A^c)\mathbb{P}(A^c)}\\
        &= \frac{0,99 \times 10^{-4}}{0,99 \times 10^{-4} + 0,001 \times 0,9999}\\
        &\approx 0,09.
\end{align*} On remarque que contrairement à l'intuition, cette
probabilité est faible.

\runinhead{Démonstration}\phantomsection\addcontentsline{toc}{subsubsection}{Démonstration}\label{I-answer-indep}
Pour \((A^c,B)\), si \(\mathbb{P}(B) = 0\), alors
\(\mathbb{P}(A^c \cap B) = 0 = \mathbb{P}(B)\mathbb{P}(A^c)\). Sinon,
\(\mathbb{P}(A^c \cap B) = \mathbb{P}(B)\mathbb{P}(A^c|B) = \mathbb{P}(B)(1-\mathbb{P}(A|B)) = \mathbb{P}(B)(1-\mathbb{P}(A))=\mathbb{P}(B)\mathbb{P}(A^c)\).

Pour \((A,B^c)\), on inverse les rôles de \(A\) et \(B\) ci-dessus.

Pour \((A^c,B^c)\), on remarque que \(A^c \cap B^c = (A \cup B)^c\),
d'où
\(\mathbb{P}(A^c \cap B^c) = 1-\mathbb{P}(A \cup B) = 1-\mathbb{P}(A) - \mathbb{P}(B) +\mathbb{P}(A)\mathbb{P}(B) = (1-\mathbb{P}(A))(1-\mathbb{P}(B)) = \mathbb{P}(A^c)\mathbb{P}(B^c)\)

\runinhead{Auto-indépendant\,?}\phantomsection\addcontentsline{toc}{subsubsection}{Auto-indépendant\,?}\label{I-answer-autoindep}
\(A\) est indépendant de lui-même ssi
\(\mathbb{P}(A) = \mathbb{P}(A \cap A) = \mathbb{P}(A)\mathbb{P}(A)\),
autrement dit si \(\mathbb{P}(A) = 0\) ou si \(\mathbb{P}(A) = 1\).

\runinhead{Durée de
fonctionnement}\phantomsection\addcontentsline{toc}{subsubsection}{Durée
de fonctionnement}\label{I-answer-ex.expo}

\begin{enumerate}
\def\labelenumi{\arabic{enumi}.}
\tightlist
\item
  \(\mathbb{P}(X \in [50,150]) = \int_{50}^{150} \frac{1}{100}\exp\left(-\frac{x}{100}\right) dx = \exp(-1/2)-\exp(-3/2) \approx 0,38.\)
\item
  \(\mathbb{P}(X \leq 100) = \int_{0}^{100} \frac{1}{100}\exp\left(-\frac{x}{100}\right) dx = 1-e^{-1} \approx 0,63.\)
\item
  \(\mathbb{P}(X = 100) = 0\)
\end{enumerate}

\runinhead{Précipitations}\phantomsection\addcontentsline{toc}{subsubsection}{Précipitations}\label{I-answer-ex.pluie}

\begin{enumerate}
\def\labelenumi{\arabic{enumi}.}
\item
  \[F(x) = \left\{ \begin{array}{ll}
  0 & \text{si } x < 0 \\
  1/2 & \text{si } x = 0 \\ 
  1/2 + 1/2 \int_0^x \frac{1}{10}\exp\left(-\frac{x}{10}\right) & \text{si } x > 0 
  \end{array} \right.\]
\item
  \(F\) n'est pas continue, \(X\) n'admet donc pas de densité
  (\protect\hyperref[I-fdr-va-dens]{voir la
  proposition (p.~\pageref*{I-fdr-va-dens})}).
\end{enumerate}

\section{Indépendance et
conditionnement}\label{I-induxe9pendance-et-conditionnement-1}

\runinhead{Question
1}\phantomsection\addcontentsline{toc}{subsubsection}{Question
1}\label{I-answer-ic-1}
On a
\(1 = \mathbb{P}(A\cup B) = \mathbb{P}(A) + \mathbb{P}(B) - \mathbb{P}(A\cap B)\),
donc \(1 = \mathbb{P}(A) + \mathbb{P}(B) - \mathbb{P}(A)\mathbb{P}(B)\)
si et seulement si \(A\) et \(B\) sont indépendants. Cette égalité se
réécrit \((1-\mathbb{P}(A))(1-\mathbb{P}(B))=0\), qui est vérifiée ssi
\(A\) ou \(B\) est presque sûr.

\runinhead{Question
2}\phantomsection\addcontentsline{toc}{subsubsection}{Question
2}\label{I-answer-ic-2}
Par définition, on a \(A\) et \(B\) indépendants relativement à
\(\mathbb{P}(\,\,|C)\) ssi
\[\mathbb{P}(A\cap B |C) = \mathbb{P}(A|C)\mathbb{P}(B|C) = \mathbb{P}(A)\mathbb{P}(B),\]
car \(A\) et \(B\) sont chacun indépendants de \(C\), d'où
\[\frac{\mathbb{P}(A\cap B \cap C)}{\mathbb{P}(C)} = \mathbb{P}(A)\mathbb{P}(B)\]
par
\protect\hyperref[I-defprobacond]{définition (p.~\pageref*{I-defprobacond})}\,;
et comme \(A\), \(B\) et \(C\) sont indépendants 2 à 2,
\[\mathbb{P}(A\cap B \cap C) = \mathbb{P}(A)\mathbb{P}(B)\mathbb{P}(C)\]
équivaut à \(A\), \(B\) et \(C\) sont mutuellement indépendants.

\runinhead{Question
3}\phantomsection\addcontentsline{toc}{subsubsection}{Question
3}\label{I-answer-ic-3}
On a
\(\mathbb{P}(A^c\cap B^c) = \mathbb{P}((A \cup B)^c) = 1 - \mathbb{P}(A) - \mathbb{P}(B) + \mathbb{P}(A\cap B)\),
donc
\[\mathbb{P}(A\cap B)\mathbb{P}(A^c \cap B^c) = \mathbb{P}(A\cap B) - \mathbb{P}(A\cap B)(\mathbb{P}(A) + \mathbb{P}(B) - \mathbb{P}(A\cap B))\]
Par ailleurs, on a \[\left\{ \begin{array}{ll}
        \mathbb{P}(A \cap B^c)  & = \mathbb{P}(A \setminus (A\cap B)) = \mathbb{P}(A) - \mathbb{P}(A\cap B) \\
        \mathbb{P}(A^c \cap B)  &=  \mathbb{P}(B \setminus (A\cap B)) = \mathbb{P}(B) - \mathbb{P}(A\cap B)
        \end{array}
        \right.\] Donc
\[ \mathbb{P}(A \cap B^c)\mathbb{P}(A^c\cap B) = \mathbb{P}(A)\mathbb{P}(B) -\mathbb{P}(A\cap B)(\mathbb{P}(A) + \mathbb{P}(B)- \mathbb{P}(A\cap B))\]
On a donc bien l'égalité ssi
\(\mathbb{P}(A)\mathbb{P}(B) = \mathbb{P}(A\cap B)\)

\section{Problème de Monty Hall}\label{I-answer-mh}

Considérons les événements \(A\) = `la porte choisie en premier est la
bonne' et \(B\) = `la dernière porte est la bonne'.

On a évidemment \(\mathbb{P}(A) = \frac{1}{3}\), puisque le choix
initial se fait uniformément parmi les 3 portes.

Pour \(B\), on doit tenir compte de l'information donnée par le
présentateur. On a par \protect\hyperref[I-formprobatot]{la formule des
probabilités totales (p.~\pageref*{I-formprobatot})}
\[\mathbb{P}(B) = \mathbb{P}(B|A)\mathbb{P}(A) + \mathbb{P}(B|A^c)\mathbb{P}(A^c) = 0.\frac{1}{3} + 1.\frac{2}{3} = \frac{2}{3}\]
En effet \(\mathbb{P}(B|A^c) = 1\) car si la porte choisie initialement
n'est pas la bonne, c'est nécessairement la dernière.

Il convient donc de changer son choix compte tenu de la nouvelle
information. Pour se convaincre du résultat, on peut refaire le calcul
avec disons 100 portes et le présentateur qui ouvre 98 autres portes
après le choix du candidat.

\section{Mesure de Probabilité}\label{I-mesure-de-probabilituxe9-1}

\runinhead{Question
1}\phantomsection\addcontentsline{toc}{subsubsection}{Question
1}\label{I-answer-mes-proba1}
On vérifie facilement les axiomes~:

\begin{itemize}
\tightlist
\item
  dans \([0,1]\)
\item
  \(Q(\Omega)= 1\)
\item
  la \(\sigma\)-additivité s'obtient en décomposant l'intégrale sur des
  \(A_n\) disjoints
\end{itemize}

\runinhead{Question
2}\phantomsection\addcontentsline{toc}{subsubsection}{Question
2}\label{I-answer-mes-proba2}
Si \(\mathbb{P}(A)=0\), alors
\(Q(A) = \int_{\Omega} X(\omega) 1_A(\omega)\mathbb{P}(d\omega) = \int_{A} X(\omega)\mathbb{P}(d\omega) = 0\)
car on intègre \(X\) sur \(A\) \(\mathbb{P}\)-négligeable.

contre-exemple~:
\((\Omega,\mathcal{A},\mathbb{P}) = (\mathbb{R},\mathcal{B}(\mathbb{R}),\mathbb{P})\)
avec \(\mathbb{P}\) symétrique et qui charge tout \(\mathbb{R}\) (la loi
normale centrée réduite par exemple) et \(X \sim \mathcal{B}_{1/2}\)
(loi de Bernoulli), définie telle que \(X(\omega) = 0\) si
\(\omega \in ]-\infty,0] = A\). On a alors \(\mathbb{P}(A) = 1/2\),
\(Q(A) = 0\) et on vérifie bien les hypothèses.

\runinhead{Question
3}\phantomsection\addcontentsline{toc}{subsubsection}{Question
3}\label{I-answer-mes-proba3}
\(Y\) positive ou bornée donc \(\mathbb{E}_Q(Y)\) est bien définie.

\(\mathbb{E}_Q(Y) = \int_{\Omega} Y(\omega) Q(d\omega) = \int_{\Omega} Y(\omega) X(\omega) \mathbb{P}(d\omega) = \mathbb{E}_\mathbb{P}(XY)\)

\runinhead{Question
4}\phantomsection\addcontentsline{toc}{subsubsection}{Question
4}\label{I-answer-mes-proba4}

\begin{enumerate}
\def\labelenumi{\roman{enumi})}
\item
  \(\mathbb{P}(X > 0) = 1\), donc \(\exists N \in \mathcal{A}\) tq
  \(\mathbb{P}(N) = 0\) et \(\forall \, \omega \in N^c,~ X(\omega) > 0\)

  Par conséquent, d'après la q.3

  \[\mathbb{E}_Q\left(\frac{1}{X}\right) = \int_{N^c} \frac{1}{X(\omega)}X(\omega) \mathbb{P}(d\omega) = 1\]
\item
  Soit \(A \in \mathcal{A}\), on a
  \(R(A) = \mathbb{E}_Q(\frac{1}{X}1_A) = \int_{A \cap N^c} \frac{1}{X(\omega)}X(\omega) \mathbb{P}(d\omega) = \mathbb{P}(A)\)
\item
  On a \(\forall \omega \in N^c, ~ X(\omega) > 0\) et donc
  \(1_{A \cap N^c}(\omega)\,X(\omega) \geq 0\). Par conséquent, si
  \(Q(A\cap N^c) = 0\), alors
  \(\mathbb{E}_Q(A\cap N^c) = \int_{A \cap N^c} X(\omega) \mathbb{P}(d\omega) = 0\)
  ce qui n'a lieu que si \(\mathbb{P}(A\cap N^c) = 0\).
\end{enumerate}

Enfin \(A \cap N \subset N\) et \(\mathbb{P}(N) = 0\) d'où
\(\mathbb{P}(A \cap N) = 0\)

\section{Queues de distributions}\label{I-queues-de-distributions-1}

\runinhead{Etude en
\(+\infty\)}\phantomsection\addcontentsline{toc}{subsubsection}{Etude en
\(+\infty\)}\label{I-answer-dlfdr-pi}

\begin{enumerate}
\def\labelenumi{\arabic{enumi}.}
\tightlist
\item
  Soit \(u>0\). Comme \(h\) est intégrable, on peut écrire
  \begin{align*}
  \int_{-\infty}^{+\infty} h(x)\,dx = \int_{-\infty}^{+\infty} x\,f(x)\,dx & = \int_{-\infty}^{u} x\,f(x)\,dx + \int_{u}^{+\infty} x\,f(x)\,dx\\
  &\geq \int_{-\infty}^{u} x\,f(x)\,dx + u\,\int_{u}^{+\infty} f(x)\,dx\\
  &= \int_{-\infty}^{u} h(x)\,dx + u\,\left(1- F(u)\right).
  \end{align*}
\end{enumerate}

\begin{enumerate}
\def\labelenumi{\arabic{enumi}.}
\setcounter{enumi}{1}
\tightlist
\item
  Soit toujours \(u>0\). Toute fonction de répartition étant à valeurs
  dans \([0,1]\), on a \(u\,\left(1-F(u)\right) \geq 0\). En utilisant
  la question précédente on obtient donc l'encadrement \begin{align*}
  0 \leq u\,\left(1-F(u)\right) & \leq \int_{-\infty}^{+\infty} h(x)\,dx - \int_{-\infty}^u h(x)\,dx\\
  &= \int_{u}^{+\infty} h(x)\,dx = \int_{-\infty}^{+\infty} h(x)\,1_{[u,+\infty]}(x)\,dx.
  \end{align*} Or pour tout \(x\in \mathbb{R}\) on a
  \(\left|h(x)\,1_{[u,+\infty]}(x)\right| \leq \left|h(x)\right|\), qui
  est intégrable. On peut donc appliquer le
  \href{Calcul\%20Intégral\%20II.pdf\%20\#TCD}{théorème de convergence
  dominée}, qui nous donne~: \begin{align*}
  \lim_{u\to+\infty} u\,(1-F(u)) & = \int_{-\infty}^{+\infty} h(x)\,\lim_{u\to+\infty} 1_{[u,+\infty]}(x)\,dx\\
  & = \int_{-\infty}^{+\infty} h(x)\,1_{\varnothing}(x)\,dx = 0.
  \end{align*} On en conclut que \(1-F(x) = o\left(\dfrac{1}{x}\right)\)
  lorsque \(x\to+\infty\).
\end{enumerate}

\runinhead{Etude en
\(-\infty\)}\phantomsection\addcontentsline{toc}{subsubsection}{Etude en
\(-\infty\)}\label{I-answer-dlfdr-mi}

\begin{enumerate}
\def\labelenumi{\arabic{enumi}.}
\setcounter{enumi}{2}
\tightlist
\item
  Prenons maintenant \(u<0\). Comme \(h\) est intégrable, on peut écrire
  \begin{align*}
  \int_{-\infty}^{+\infty} h(x)\,dx = \int_{-\infty}^{+\infty} x\,f(x)\,dx & = \int_{-\infty}^{u} x\,f(x)\,dx + \int_{u}^{+\infty} x\,f(x)\,dx\\
  &\leq u\, \int_{-\infty}^{u} f(x)\,dx + \int_{u}^{+\infty} x\,f(x)\,dx\\
  &= u\,F(u) + \int_{u}^{+\infty} h(x)\,dx.
  \end{align*} Or, comme \(F\) prend ses valeurs dans \([0,1]\),
  \(u\,F(u) \leq 0\). On obtient donc \begin{align*}
  0 \geq u\,F(u) & \geq \int_{-\infty}^{+\infty} h(x)\,dx - \int_{u}^{+\infty} h(x)\,dx\\
  &= \int_{-\infty}^{u} h(x)\,dx = \int_{-\infty}^{+\infty} h(x)\,1_{[-\infty,u]}(x)\,dx.
  \end{align*} Comme précédemment, cette dernière fonction sous
  l'intégrale est dominée par \(-\left|h\right|\) (à gauche) et
  \(\left|h\right|\) (à droite), toutes deux intégrables, donc d'après
  le \href{Calcul\%20Intégral\%20II.pdf\%20\#TCD}{théorème de
  convergence dominée} \begin{align*}
  \lim_{u\to-\infty} u\,F(u) & = \int_{-\infty}^{+\infty} h(x)\,\lim_{u\to-\infty} 1_{[-\infty,u]}(x)\,dx\\
  & = \int_{-\infty}^{+\infty} h(x)\,1_{\varnothing}(x)\,dx = 0.
  \end{align*} On a donc bien \(F(x) = o\left(\dfrac{1}{x}\right)\)
  lorsque \(x\to-\infty\).
\end{enumerate}

\runinhead{Interprétation}\phantomsection\addcontentsline{toc}{subsubsection}{Interprétation}\label{I-answer-dlfdr-interpret}

\begin{enumerate}
\def\labelenumi{\arabic{enumi}.}
\setcounter{enumi}{3}
\tightlist
\item
  L'hypothèse que \(h\) est intégrable nous donne une indication sur la
  vitesse de convergence de \(1-F(x)\) et \(F(x)\) lorsque \(x\) tend
  respectivement vers \(+\infty\) et \(-\infty\)~: elles convergent vers
  \(0\) au moins aussi vite que \(\dfrac{1}{x}\). On dit que les queues
  de la distribution sont ``relativement'' fines.
\end{enumerate}

Pour mieux comprendre l'impact de ce résultat, prenons un exemple
pratique. Considérons un phénomène aléatoire comme la concentration en
polluant dans l'air, modélisé par une loi de probabilité sur
\(\mathbb{R}\) de densité \(f\) nulle sur \(\mathbb{R}_-^\ast\). Etre
capable de déterminer la probabilité que la concentration dépasse un
seuil critique \(x\) est alors très important pour les organismes de
contrôle de la qualité de l'air. Or cet événement correspond
mathématiquement à l'événement \([x,+\infty]\). Dans ce cas,
\(1 - F(x)\) nous donne cette probabilité selon le modèle considéré. Le
résultat nous indique que si l'on choisit un modèle tel que \(h\) est
intégrable, alors cette probabilité décroît plus rapidement que
\(\dfrac{1}{x}\) lorsque le seuil \(x\) augmente. En d'autres termes,
les événements extrêmes (quand \(x\) est grand) restent ``assez'' rares.
L'impact du choix de modèle n'est donc pas négligeable~: si on prend une
distribution à queues trop fines alors que les pics de pollution ne sont
en réalité pas si rares que ça, \(1-F(x)\) risque de sous-estimer le
risque réel de dépassement du seuil\,!

\runinhead{Etude de la
réciproque}\phantomsection\addcontentsline{toc}{subsubsection}{Etude de
la réciproque}\label{I-answer-dlfdr-reciproque}

\begin{enumerate}
\def\labelenumi{\arabic{enumi}.}
\setcounter{enumi}{4}
\item
  La fonction \(F\) considérée est continue et dérivable par morceaux
  sur \(\mathbb{R}\)\,; le seul point où \(F\) n'est pas dérivable est
  \(e\) (le taux d'accroissement n'a pas les mêmes limites à gauche et à
  droite). Elle possède donc \protect\hyperref[I-fdr-va-dens]{une
  densité (p.~\pageref*{I-fdr-va-dens})} qui s'obtient en dérivant
  chaque morceau~: pour tout \(x\in\mathbb{R}\),
  \[f(x) = \left|\begin{array}{ll} e\,\dfrac{\ln(x) + 1}{x^2\,\ln(x)^2} & \text{si } x \geq e,\\ 0 & \text{sinon.}\end{array}\right.\]
\item
  On remarque que pour tout \(x\geq e\),
  \(x\,\left(1-F(x)\right) = \dfrac{e}{\ln(x)} \xrightarrow[x\to+\infty]{} 0\),
  donc le développement asymptotique \((1)\) est bien vérifié. En outre,
  comme \(x\,F(x) = 0\) pour tout \(x < e\), le développement
  asymptotique \((2)\) est aussi bien respecté.
\item
  Pour tout \(x\in\mathbb{R}\),
  \[ h(x) = \left|h(x)\right| = \left|\begin{array}{ll} e\,\dfrac{\ln(x) + 1}{x\,\ln(x)^2} & \text{si } x \geq e,\\ 0 & \text{sinon,} \end{array}\right.\]
  et on pose \(h(\pm\infty) = \lim\limits_{x\to\pm\infty}h(x) = 0\). Par
  conséquent, \(h\) est clairement intégrable sur \([-\infty,e]\), où
  elle est égale presque partout (partout sauf en \(e\)) à la fonction
  nulle. Pour savoir si \(h\) est intégrable sur toute la droite réelle
  il nous faut donc regarder si elle l'est sur \([e,+\infty]\). On
  remarque en outre que pour tout \(x\geq e\),
  \[h(x) \geq \dfrac{1}{x\,\ln(x)}, \] et le terme de droite n'est pas
  intégrable (en effet, une primitive de \(\frac1{x \ln(x)}\) est
  \(\ln(ln(x))\) qui tend vers \(+\infty\) en \(+\infty\)). Par
  conséquent, \(h\) n'est pas intégrable.
\item
  La fonction de répartition \(F\) considérée respecte bien les
  développements asymptotiques \((1)\) et \((2)\), mais la fonction
  \(h\) associée n'est pas intégrable. Nous avons donc montré que la
  réciproque de la propriété est fausse.
\end{enumerate}

\runinhead{Liens entre les développements
asymptotiques}\phantomsection\addcontentsline{toc}{subsubsection}{Liens
entre les développements asymptotiques}\label{I-answer-dlfdr-liensdl}

\begin{enumerate}
\def\labelenumi{\arabic{enumi}.}
\setcounter{enumi}{8}
\tightlist
\item
  On remarque que \(F\) satisfait \((2)\) mais pas \((1)\), tandis que
  \(G\) satisfait \((1)\) mais pas \((2)\). On en conclut que \((1)\)
  n'implique pas \((2)\) et inversement.
\end{enumerate}

\section{Densité et fonction de répartition d'une loi
Normale}\label{I-densituxe9-et-fonction-de-ruxe9partition-dune-loi-normale-1}

\runinhead{Propriétés
générales}\phantomsection\addcontentsline{toc}{subsubsection}{Propriétés
générales}\label{I-answer-fdrgaussprop}

\begin{enumerate}
\def\labelenumi{\arabic{enumi}.}
\tightlist
\item
  Faisons l'étude de \(f\).
\end{enumerate}

\begin{itemize}
\tightlist
\item
  \textbf{Domaine.} \(f : \mathbb{R}\to \mathbb{R}^+\).
\item
  \textbf{Parité.} On remarque que \(f\) est paire~:
  \(\forall\,x\in\mathbb{R}\), \(f(x) = f(-x)\).
\item
  \textbf{Limites aux bornes.}
  \(\lim\limits_{x\to-\infty} f(x) = \lim\limits_{x\to+\infty} f(x) = 0\).
  Ainsi, la parité est aussi vraie aux bornes.
\item
  \textbf{Dérivabilité.} \(f\) est continue sur \(\mathbb{R}\) et
  infiniment dérivable sur \(\mathbb{R}\) en tant que composée de
  fonctions infiniment dérivables. Pour tout \(x\in\mathbb{R}\) sa
  dérivée première s'écrit \(f^\prime(x) = -x\,f(x)\) et a dérivée
  seconde \(f^{\prime\prime}(x) =(x^2 - 1)\,f(x)\).
\item
  \textbf{Variations.} \(f^\prime\) est strictement positive (resp.
  négative) ssi \(-\infty < x < 0\) (resp. \(+\infty > x > 0\)). Elle
  est nulle en \(0\), \(+\infty\) et \(-\infty\). Ainsi, \(f\) est
  strictement croissante sur \(]-\infty,0[\), vaut
  \(\left(2\,\pi\right)^{-1/2}\) en \(0\), puis est strictement
  décroissante sur \(]0,+\infty[\).
\item
  \textbf{Convexité/Concavité.} \(f^{\prime\prime}\) est strictement
  positive (resp. négative) sur \(\mathbb{R}\backslash\,]-1,1[\) (resp.
  \(]-1,1[\)). \(f\) est donc convexe sur
  \(\mathbb{R}\backslash\,]-1,1[\) et concave sur \(]-1,1[\).
\end{itemize}

\begin{enumerate}
\def\labelenumi{\arabic{enumi}.}
\setcounter{enumi}{1}
\tightlist
\item
  Par définition, pour tout \(x\in\mathbb{R}\)
  \[ F(x) = \int_{-\infty}^x \dfrac{1}{\sqrt{2\,\pi}}\,\exp\left\{-\dfrac{u^2}{2}\right\}\,du.\]
  On pose \(F(+\infty) = 1\) et \(F(-\infty) = 0\). Lorsque \(x = 0\),
  comme \(f\) est paire et que son intégrale sur \([-\infty,+\infty]\)
  vaut \(1\) (c'est \protect\hyperref[I-defdens]{une
  densité (p.~\pageref*{I-defdens})}), son intégrale sur \([-\infty,0]\)
  (i.e.~\(F(0)\)) vaut \(\dfrac{1}{2}\). Lorsque \(x > 0\), en
  décomposant l'intégrale puis en utilisant un
  \href{Calcul\%20Intégral\%20I.pdf\%20\#changement-de-variable}{changement
  de variable} on obtient
  \[F(x) = \dfrac{1}{2} + \int_{0}^x \dfrac{1}{\sqrt{2\,\pi}}\,\exp\left\{-\dfrac{u^2}{2}\right\}\,du = \dfrac{1}{2} + \int_{0}^{\frac{x}{\sqrt{2}}} \dfrac{1}{\sqrt{\pi}}\,e^{-v^2}\,dv = \dfrac{1}{2} + \dfrac{1}{2}\,\text{erf}\left(\dfrac{x}{\sqrt{2}}\right).\]
  Lorsque \(x<0\), on peut procéder de la même manière avec \(1 - F(x)\)
  et obtenir la même égalité que ci-dessus.
\end{enumerate}

\runinhead{Encadrement de \(1 - F(x)\) pour tout
\(x > 0\)}\phantomsection\addcontentsline{toc}{subsubsection}{Encadrement
de \(1 - F(x)\) pour tout \(x > 0\)}\label{I-answer-fdrgaussenca}

\begin{enumerate}
\def\labelenumi{\arabic{enumi}.}
\setcounter{enumi}{2}
\item
  Soit \(x>0\). Par définition,
  \[1 - F(x) = \int_{-\infty}^{+\infty} f(u)\,du - \int_{-\infty}^x f(u)\,du = \int_{x}^{+\infty} \dfrac{1}{\sqrt{2\,\pi}}\,\exp\left\{-\dfrac{u^2}{2}\right\}\,du. \]
  Or pour tout \(u\geq x\) on a \(u/x \geq 1\), donc
  \[1 - F(x) \leq \dfrac{1}{\sqrt{2\,\pi}}\,\int_{x}^{+\infty} \dfrac{u}{x}\, \exp\left\{-\dfrac{u^2}{2}\right\}\,du = \dfrac{-1}{x}\,\int_{x}^{+\infty} f^\prime(u)\,du. \]
  Or \(f^\prime\) a pour primitive \(f\) donc on a
  \[1-F(x) \leq -\dfrac{f(+\infty) - f(x)}{x} = \dfrac{f(x)}{x}. \]
\item
  Soit \(x>0\). Par définition,
  \[\left(1 + \dfrac{1}{x^2}\right)\,\left(1-F(x)\right) = \int_x^{+\infty} \left(1 + \dfrac{1}{x^2} \right)\,f(u)\,du.\]
  Or pout tout \(u \geq x\) on a \(\dfrac{1}{u^2} \leq \dfrac{1}{x^2}\),
  d'où
  \[\left(1 + \dfrac{1}{x^2}\right)\,\left(1-F(x)\right) \geq \int_{x}^{+\infty} \left( 1 + \dfrac{1}{u^2} \right)\,f(u)\,du.\]
  Remarquons maintenant que d'après la question 1,
  \(\left(\dfrac{f(x)}{x} \right)^\prime = \dfrac{-x^2\,f(x) - f(x)}{x^2} = -\left(1+\dfrac{1}{x^2}\right)\,f(x)\).
  En d'autres termes, \(A : x\in\mathbb{R}\mapsto -\dfrac{f(x)}{x}\) est
  une primitive de
  \(x\in\mathbb{R}\mapsto \left(1+\dfrac{1}{x^2}\right)\,f(x)\). On
  obtient alors par
  \[\left(1 + \dfrac{1}{x^2}\right)\,\left(1-F(x)\right) \geq - \left(A(+\infty) - A(x)\right) = A(x) = \dfrac{f(x)}{x}.\]
  On en déduit directement l'inégalité de gauche~:
  \[ 1-F(x) \geq f(x)\,\dfrac{x^2}{(x^2+1)\,x} = f(x)\,\dfrac{x}{x^2+1}. \]
\end{enumerate}

\runinhead{Equivalent de \(1-F(x)\) quand
\(x\to+\infty\).}\phantomsection\addcontentsline{toc}{subsubsection}{Equivalent
de \(1-F(x)\) quand \(x\to+\infty\).}\label{I-answer-fdrgaussequi}

\begin{enumerate}
\def\labelenumi{\arabic{enumi}.}
\setcounter{enumi}{4}
\item
  Soit \(x>0\). L'encadrement que nous avons démontré peut être réécrit
  \[ \dfrac{x^2}{1+x^2} \leq \dfrac{1-F(x)}{\dfrac{f(x)}{x}} \leq 1, \]
  et comme \(\dfrac{x^2}{1+x^2} \to 1\) lorsque \(x\to+\infty\), on
  obtient que \(1-F(x) \sim \dfrac{f(x)}{x}\) quand \(x\to+\infty\).
\item
  Pour tout \(x \in [-\infty,+\infty]\) on a
  \(\left|h(x)\right| = \left|x\,f(x)\right| = |x|\,f(x)\) puisque
  \(f\geq 0\). Lorsque \(x \in [0,+\infty]\),
  \(\left|h(x)\right| = h(x)\), or d'après la question 1, \(h\) admet
  pour primitive \(-f\), qui vaut \(0\) en \(+\infty\) et \(-\infty\)
  puis \((2\,\pi)^{-1/2}\) en \(0\). Par conséquent, (d'après
  l'extension du \href{Calcul\%20Intégral\%20I.pdf\%20\#TFC}{théorème
  fondamental du calcul}) l'intégrale de \(\left|h\right|\) sur
  \([0,+\infty]\) existe, et vaut \((2\,\pi)^{-1/2}\). Or
  \(\left|h\right|\) est paire, son intégrale sur \([-\infty,0]\) existe
  donc aussi et vaut \((2\,\pi)^{-1/2}\). On en conclut que \(h\) est
  bien absolument intégrable sur \([-\infty,+\infty]\). En remarquant
  qu'elle est impaire, on obtient immédiatement que son intégrale vaut
  \(0\).
\end{enumerate}

On est donc bien dans le cadre de l'exercice sur les queues de
distributions. L'équivalent que nous avons trouvé implique bien que
\(1-F(x) = o\left(\dfrac{1}{x} \right)\) lorsque \(x\to+\infty\). Il
nous donne juste plus de précisions quant au comportement de \(1-F(x)\)
lorsque \(x\to+\infty\) dans le cas spécifique de la loi Normale centrée
réduite~: on connaît sa vitesse de convergence.

\runinhead{Loi Normale
générale.}\phantomsection\addcontentsline{toc}{subsubsection}{Loi
Normale générale.}\label{I-answer-fdrgaussgen}

\begin{enumerate}
\def\labelenumi{\arabic{enumi}.}
\setcounter{enumi}{6}
\tightlist
\item
  Soit \(x\in\mathbb{R}\). On remarque que
  \(g(x) = \dfrac{1}{\sigma}\,f\left(\dfrac{x-\mu}{\sigma} \right)\).
  Cela implique en particulier que
  \[G(x) =  \int_{-\infty}^x \dfrac{1}{\sigma}\,f\left(\dfrac{u-\mu}{\sigma} \right)\,du = \int_{-\infty}^{\frac{x-\mu}{\sigma}} f(u)\,du = F\left(\dfrac{x-\mu}{\sigma}\right).\]
  On obtient alors l'encadrement suivant, pour tout \(x>\mu\)~:
  \[ \dfrac{x-\mu}{1 + \dfrac{(x-\mu)^2}{\sigma^2}}\,g(x) \leq 1 - G(x) \leq \dfrac{\sigma^2\,g(x)}{x-\mu}.  \]
\end{enumerate}

On en déduit que
\(1-G(x) \sim \dfrac{\sigma^2\,g(x)}{x-\mu} \sim \sigma^2\,\dfrac{g(x)}{x}\)
lorsque \(x\to+\infty\).

\section{Loi de vie et de mort}\label{I-loi-de-vie-et-de-mort-1}

\runinhead{Question
1}\phantomsection\addcontentsline{toc}{subsubsection}{Question
1}\label{I-answer-viemort1}
Par définition, on a
\[ G_{t_0}(t) = \mathbb{P}(T>t+t_0|T>t_0) = \frac{\mathbb{P}(T>t+t_0,T>t)}{\mathbb{P}(T>t)} = \frac{G(t+t_0)}{G(t)}\]

\runinhead{Question
2}\phantomsection\addcontentsline{toc}{subsubsection}{Question
2}\label{I-answer-viemort2}
D'après la question précédente, \(G\) vérifie alors
\(G(t_0+t) = G(t)G(t_0)\) pour tous \(t, t_0 > 0\). Comme \(G\) est
décroissante, continue à droite et tend vers 0 à l'infini, on en déduit
que \(G(t) = e ^{-\theta t}\), pour un \(\theta > 0\). On reconnaît la
fonction de répartition complémentaire d'une loi exponentielle de
paramètre \(\theta\).

\runinhead{Question
3}\phantomsection\addcontentsline{toc}{subsubsection}{Question
3}\label{I-answer-viemort3}

\begin{align*}
D(t) &= \lim_{\Delta t \to 0} \frac{1}{\Delta t}\mathbb{P}(t < T \leq t + \Delta t | T >t)\\
     &= \lim_{\Delta t \to 0} \frac{1}{\Delta t}\frac{\mathbb{P}(t < T \leq t + \Delta t)}{\mathbb{P}(T >t)}\\
     &= \frac{1}{G(t)} \lim_{\Delta t \to 0} \frac{G(t) - G(t + \Delta t)}{\Delta t}\\
     &= -\frac{g(t)}{G(t)}
\end{align*} Ainsi, \(D(t) = \frac{d}{dt}(-\ln G(t))\) et comme
\(G(0) = 1\), alors on a pour \(t >0\)
\[G(t) = \exp\left(-\int_0^tD(s) ds\right)\] Si \(D\) est constant, on
retrouve une loi exponentielle.
